%%% FIELD proof-aipw-martingale
Fix \(v\in\Delta(\mathcal B)\).  Put \(u_b=q_b\mathbf 1\{b\in W\}\), and, for an arm \(z\in A_b\), define
\[
e_{bi}(z)=Y_{bi}(z)-m_{bi}(z),
\qquad
\Delta_{bi}(z)=\phi_{bi}(z)-Y_{bi}(z).
\]
Consistency of the observed outcome and Definition \(\mathrm{def:aipw\text{-}score}\) give
\[
\Delta_{bi}(z)
=e_{bi}(z)\left\{\frac{\mathbf 1\{Z_{bi}=z\}}
 {p_{bi}(z\mid H_{bi})}-1\right\}.
\]
The denominator is positive by Assumption \(\mathrm{ass:availability\text{-}support}\).  Both \(m_{bi}(z)\) and \(p_{bi}(z\mid H_{bi})\) are \(H_{bi}\)-measurable.  Hence Assumption \(\mathrm{ass:sequential\text{-}randomization}\) implies, conditionally on \(H_{bi}\),
\[
\mathbb E_{\mathcal D,Y}[\Delta_{bi}(z)\mid H_{bi}]=0
\]
and
\[
\begin{aligned}
\mathbb E_{\mathcal D,Y}[\Delta_{bi}(z)^2\mid H_{bi}]
&=e_{bi}(z)^2\left[
p_{bi}(z\mid H_{bi})
 \left\{\frac1{p_{bi}(z\mid H_{bi})}-1\right\}^{2}
+1-p_{bi}(z\mid H_{bi})\right]\\
&=e_{bi}(z)^2
 \frac{1-p_{bi}(z\mid H_{bi})}{p_{bi}(z\mid H_{bi})}\\
&\le 4M^2
 \frac{1-p_{bi}(z\mid H_{bi})}{p_{bi}(z\mid H_{bi})}.
\end{aligned}
\]
The last inequality follows from Assumption \(\mathrm{ass:bounded\text{-}outcomes}\) and the stipulated bound \(m_{bi}(z)\in[-M,M]\).  Define the increment, with its first term interpreted as zero when \(b\notin W\), by
\[
\xi_{bi}(v)
=\frac{u_b}{n_b}\Delta_{bi}(z_1)
-\frac{v_b}{n_b}\Delta_{bi}(z_0).
\]
Assumption \(\mathrm{ass:availability\text{-}window}\) makes every score occurring with nonzero coefficient available.  The preceding conditional-mean identity shows that the lexicographically ordered array 
\(\{\xi_{bi}(v),H_{bi}\}_{(b,i)\in\mathcal T}\) is a martingale-difference array.  Moreover, \((x-y)^2\le 2x^2+2y^2\) and the preceding second-moment calculation yield
\[
\begin{aligned}
\mathbb E_{\mathcal D,Y}[\xi_{bi}(v)^2\mid H_{bi}]
&\le 8M^2\left\{
\frac{u_b^2}{n_b^2}
\frac{1-p_{bi}(z_1\mid H_{bi})}{p_{bi}(z_1\mid H_{bi})}
+\frac{v_b^2}{n_b^2}
\frac{1-p_{bi}(z_0\mid H_{bi})}{p_{bi}(z_0\mid H_{bi})}
\right\},
\end{aligned}
\]
again taking the first summand to be zero off \(W\).  Summing gives exactly
\[
\sum_{b,i}\mathbb E_{\mathcal D,Y}[\xi_{bi}(v)^2\mid H_{bi}]
\le \mathsf G_{\mathcal D}(v;H).
\]

It remains to identify the deterministic center.  By Definitions \(\mathrm{def:target}\), \(\mathrm{def:concurrent\text{-}weights}\), and \(\mathrm{def:borrowing\text{-}estimator}\),
\[
\begin{aligned}
\widehat\tau(v)-\tau_W
&=\sum_{b,i}\xi_{bi}(v)
+\sum_{b\in W}\frac{q_b}{n_b}\sum_iY_{bi}(z_1)
-\sum_b\frac{v_b}{n_b}\sum_iY_{bi}(z_0)\\
&\qquad-
\sum_{b\in W}\frac{q_b}{n_b}\sum_i
 \{Y_{bi}(z_1)-Y_{bi}(z_0)\}\\
&=\sum_{b,i}\xi_{bi}(v)-\sum_b(v_b-u_b)\mu_b\\
&=\sum_{b,i}\xi_{bi}(v)-\sum_bd_b(v)\mu_b.
\end{aligned}
\]
Finally, taking expectations in the conditional-variance bound and using Definition \(\mathrm{def:average\text{-}design\text{-}scale}\) gives, for every \(Y\in\mathcal Y_{M,\gamma}\),
\[
\mathbb E_{\mathcal D,Y}\sum_{b,i}\xi_{bi}(v)^2
\le \mathbb E_{\mathcal D,Y}[\mathsf G_{\mathcal D}(v;H)]
\le\overline{\mathsf V}_{\mathcal D,\gamma}(v),
\]
which proves all assertions.

%%% FIELD proof-tv-support-dual
Fix \(v\in\Delta(\mathcal B)\) and set \(D_0(v)=0\) and \(D_B(v)=\sum_{b=1}^Bd_b(v)\).  Both \(v\) and \(u\) are probability vectors, so \(D_B(v)=0\).  Since \(d_b(v)=D_b(v)-D_{b-1}(v)\), finite summation by parts gives
\[
\begin{aligned}
\sum_{b=1}^Bd_b(v)m_b
&=\sum_{b=1}^B\{D_b(v)-D_{b-1}(v)\}m_b\\
&=\sum_{j=1}^{B-1}D_j(v)(m_j-m_{j+1})\\
&=-\sum_{j=1}^{B-1}D_j(v)(m_{j+1}-m_j).
\end{aligned}
\]
Write \(c=m_1\) and \(\delta_j=m_{j+1}-m_j\).  Then
\[
m_b=c+\sum_{k<b}\delta_k.
\]
Thus the feasible set in Definition \(\mathrm{def:drift\text{-}support\text{-}family}\) is represented exactly by
\[
\sum_{j=1}^{B-1}|\delta_j|\le\gamma,
\qquad
-M\le c+\sum_{k<b}\delta_k\le M
\quad(1\le b\le B).
\]
After replacing each \(\delta_j\) by the difference of two nonnegative variables, the absolute-value constraint is linear.  Maximizing each of the two signs of
\(-\sum_jD_j(v)\delta_j\) therefore gives two finite linear programs, and their larger value is precisely \(\mathsf B_\gamma(v)\).  The feasible polytope is nonempty and compact, so a maximizing extreme point exists.  Given any maximizing \(m\), the fixed schedule with \(Y_{bi}(z_0)=m_b\) for every \(i\), and with every other available potential outcome set to zero, is bounded by \(M\), has control means \(m_b\), and obeys the same total-variation constraint.  It is therefore an extremal fixed control schedule in \(\mathcal Y_{M,\gamma}\); no stochastic or superpopulation construction is being used.

For \(0\le\gamma\le2M\), the summation-by-parts identity and the variation constraint imply
\[
\left|\sum_bd_b(v)m_b\right|
\le \max_{j<B}|D_j(v)|\sum_{j=1}^{B-1}|\delta_j|
\le\gamma\max_{j<B}|D_j(v)|.
\]
Choose \(j_\star\) attaining the maximum.  If \(D_{j_\star}(v)\ne0\), put
\[
\delta_{j_\star}=-\gamma\operatorname{sgn}(D_{j_\star}(v)),
\qquad \delta_j=0\ (j\ne j_\star),
\qquad c=\frac{\gamma}{2}\operatorname{sgn}(D_{j_\star}(v)).
\]
The resulting block means equal \(c\) before the jump and \(-c\) after the jump, so they lie in \([-M,M]\) because \(\gamma\le2M\), and the absolute objective is
\(\gamma|D_{j_\star}(v)|\).  If every \(D_j(v)\) is zero, both sides are zero.  Hence
\[
\mathsf B_\gamma(v)=\gamma\max_{j<B}|D_j(v)|,
\qquad 0\le\gamma\le2M.
\]

Finally, every \(m\in[-M,M]^B\) has
\(\sum_{j<B}|m_{j+1}-m_j|\le2M(B-1)=\Gamma_{\mathrm{vac}}\).  The variation constraint is therefore vacuous at that budget.  Choosing \(m_b=M\operatorname{sgn}(d_b(v))\), with an arbitrary sign when \(d_b(v)=0\), attains the elementary box support function, and consequently
\[
\mathsf B_{\Gamma_{\mathrm{vac}}}(v)
=\max_{m\in[-M,M]^B}\left|\sum_bd_b(v)m_b\right|
=M\sum_b|d_b(v)|.
\]

%%% FIELD proposed-rare-statement
Fix any enrollment position \((b,i)\), schedule \(Y\), and arm \(z\) such that \(p_{bi}(z\mid H_{bi})>0\) holds \(P^O_{\mathcal D,Y}\)-almost surely, and fix any predictable event \(E\in H_{bi}\).  The event's contribution to the average inverse-propensity factor is exactly 
\(\mathbb E_{\mathcal D,Y}[\mathbf 1_E\{1-p_{bi}(z\mid H_{bi})\}/p_{bi}(z\mid H_{bi})]\).  If \(P^O_{\mathcal D,Y}(E)\le\epsilon\) and \(p_{bi}(z\mid H_{bi})\ge\delta\) on \(E\), this contribution is at most \(\epsilon(1-\delta)/\delta\).  Thus a rare \(\epsilon\)-probability branch is charged through its assignment probability rather than by the worst-path factor \(1/\delta\) alone.

%%% FIELD proof-rare-history
Let \(p=p_{bi}(z\mid H_{bi})\).  The added positivity condition makes the displayed random variable well defined.  Its contribution on \(E\) is, by restriction of the expectation under the actual design law,
\[
\mathbb E_{\mathcal D,Y}\left[\mathbf 1_E\frac{1-p}{p}\right].
\]
On \(E\), the hypothesis \(p\ge\delta>0\) and monotonicity of \(x\mapsto x^{-1}-1\) on \((0,1]\) give
\[
0\le\frac{1-p}{p}\le\frac{1-\delta}{\delta}.
\]
Therefore
\[
\mathbb E_{\mathcal D,Y}\left[\mathbf 1_E\frac{1-p}{p}\right]
\le \frac{1-\delta}{\delta}P^O_{\mathcal D,Y}(E)
\le\epsilon\frac{1-\delta}{\delta}.
\]
This is an assignment-law average: the factor contributed by the branch is multiplied by its actual probability.  Assumption \(\mathrm{ass:propensity\text{-}floor}\) guarantees the positivity premise, with \(p\ge p_{\min}\), whenever the fixed arm is available.

%%% FIELD proof-honest-upper
We first establish existence of the optimizer from the primitive design conditions.  For a fixed schedule \(Y\in\mathcal Y_{M,\gamma}\), taking expectation in the displayed definition of \(\mathsf G_{\mathcal D}\) gives
\[
\mathbb E_{\mathcal D,Y}[\mathsf G_{\mathcal D}(v;H)]
=c_Y+\sum_{b=1}^Ba_{b,Y}v_b^2,
\]
where
\[
\begin{aligned}
c_Y&=8M^2\mathbb E_{\mathcal D,Y}
 \sum_{b,i}\frac{u_b^2}{n_b^2}
 \frac{1-p_{bi}(z_1\mid H_{bi})}{p_{bi}(z_1\mid H_{bi})},\\
a_{b,Y}&=8M^2\mathbb E_{\mathcal D,Y}
 \sum_{i=1}^{n_b}\frac1{n_b^2}
 \frac{1-p_{bi}(z_0\mid H_{bi})}{p_{bi}(z_0\mid H_{bi})}.
\end{aligned}
\]
Here \(c_Y\ge0\) and \(a_{b,Y}\ge0\).  Assumptions \(\mathrm{ass:availability\text{-}window}\), \(\mathrm{ass:availability\text{-}support}\), and \(\mathrm{ass:propensity\text{-}floor}\) make every nonzero displayed denominator at least \(p_{\min}\); hence these coefficients are finite and uniformly bounded over the schedule class.  Since the square root is increasing,
\[
\sqrt{\overline{\mathsf V}_{\mathcal D,\gamma}(v)/\alpha}
=\sup_{Y\in\mathcal Y_{M,\gamma}}
 \alpha^{-1/2}\sqrt{c_Y+\sum_ba_{b,Y}v_b^2}.
\]
For each fixed \(Y\), the expression after the supremum is the Euclidean norm of
\(\alpha^{-1/2}(\sqrt{c_Y},\sqrt{a_{1,Y}}v_1,\ldots,\sqrt{a_{B,Y}}v_B)\), so it is convex and continuous in \(v\).  Its supremum over \(Y\) is convex and lower semicontinuous.  Definition \(\mathrm{def:drift\text{-}support\text{-}family}\) expresses \(\mathsf B_\gamma\) as the support function of a nonempty compact set under an affine map of \(v\); it is finite, convex, and continuous.  Consequently \(\Psi_{\alpha,\gamma}\) is finite, convex, and lower semicontinuous.  The simplex \(\Delta(\mathcal B)\) is nonempty and compact, so a minimizing sequence has a convergent subsequence and lower semicontinuity shows that its limit is an optimizer \(v_\gamma^\star\).

Now fix \(Y\in\mathcal Y_{M,\Gamma}\), put \(v=v^\star\), and write \(S=\sum_{b,i}\xi_{bi}(v)\).  For two lexicographically ordered increments with the first preceding the second, the tower property gives \(\mathbb E[\xi_s\xi_t]=\mathbb E[\xi_s\mathbb E(\xi_t\mid H_t)]=0\).  Lemma \(\mathrm{lem:aipw\text{-}martingale}\) and Definition \(\mathrm{def:average\text{-}design\text{-}scale}\) therefore imply
\[
\mathbb E_{\mathcal D,Y}[S^2]
=\sum_{b,i}\mathbb E_{\mathcal D,Y}[\xi_{bi}(v)^2]
\le\overline{\mathsf V}_{\mathcal D,\Gamma}(v).
\]
The same lemma and the definition of the sharp support functional give
\[
\widehat\tau(v)-\tau_W
=S-\sum_bd_b(v)\mu_b,
\qquad
\left|\sum_bd_b(v)\mu_b\right|\le\mathsf B_\Gamma(v),
\]
where the last inequality uses Assumption \(\mathrm{ass:control\text{-}tv}\), Assumption \(\mathrm{ass:bounded\text{-}outcomes}\), and Lemma \(\mathrm{lem:tv\text{-}support\text{-}dual}\).  If \(\overline{\mathsf V}_{\mathcal D,\Gamma}(v)>0\), nonnegativity gives
\[
\mathbb E[S^2]\ge
\frac{\overline{\mathsf V}_{\mathcal D,\Gamma}(v)}{\alpha}
P^O_{\mathcal D,Y}\left\{
S^2>\frac{\overline{\mathsf V}_{\mathcal D,\Gamma}(v)}{\alpha}
\right\}.
\]
Combining this display with the second-moment bound yields
\[
P^O_{\mathcal D,Y}\left(
 |S|>\sqrt{\overline{\mathsf V}_{\mathcal D,\Gamma}(v)/\alpha}
\right)\le\alpha.
\]
If the variance envelope is zero, the same second-moment bound gives \(S=0\) almost surely, so the corresponding conclusion holds with probability one.  Thus in all cases
\[
P^O_{\mathcal D,Y}\left(
 |\widehat\tau(v)-\tau_W|\le
 \mathsf B_\Gamma(v)+
 \sqrt{\overline{\mathsf V}_{\mathcal D,\Gamma}(v)/\alpha}
\right)\ge1-\alpha.
\]
The radius on the right is \(\Psi_\alpha(v^\star)\).

Assumption \(\mathrm{ass:bounded\text{-}outcomes}\), the positivity and unit-sum properties of \(q\), and Definition \(\mathrm{def:target}\) imply \(\tau_W\in[-2M,2M]\).  For every \(t\in[-2M,2M]\), the scalar projection satisfies
\[
|\Pi_M(x)-t|\le|x-t|;
\]
this follows directly by considering \(x<-2M\), \(x\in[-2M,2M]\), and \(x>2M\).  Therefore the preceding event implies
\(|\Pi_M(\widehat\tau(v^\star))-\tau_W|\le\Psi_\alpha(v^\star)\).  If the radius is below \(2M\), the interval centered at the projected value contains \(\tau_W\), and intersecting with \([-2M,2M]\) does not remove it.  This intersection is nonempty on every path because its center lies in \([-2M,2M]\).  If the radius is at least \(2M\), Definition \(\mathrm{def:bias\text{-}aware\text{-}interval}\) returns all of \([-2M,2M]\), which contains \(\tau_W\) surely.  Hence
\[
\inf_{Y\in\mathcal Y_{M,\Gamma}}
P^{K^\star}_{\mathcal D,Y}(\tau_W\in I^\star)\ge1-\alpha.
\]
All ingredients of \(v^\star\) are fixed by the known design and the prespecified schedule class, while \(I^\star\) is a measurable function of the observed scores.  It therefore induces the schedule-independent Dirac Markov kernel in Definition \(\mathrm{def:projected\text{-}kernel}\).

Finally, on the projected branch,
\(\operatorname{len}(I^\star)\le2\Psi_\alpha(v^\star)\); on the full-space branch its length is \(4M\).  Both cases are summarized by
\[
\operatorname{len}(I^\star)
\le2\min\{2M,\Psi_\alpha(v^\star)\}
=2\mathfrak R_\alpha^\star
\]
on every path.  Taking expectations and then the supremum over the fixed-schedule class proves the asserted worst-case expected-length bound.  The nonempty construction shows that the claimed empty-intersection fallback event has probability zero.

%%% FIELD proof-strict-improvement
Fix \(\gamma\) and abbreviate \(f(v)=\Psi_{\alpha,\gamma}(v)\) and \(C=\Delta(\mathcal B)\).  For every fixed schedule \(Y\), Definition \(\mathrm{def:average\text{-}design\text{-}scale}\) writes the expected variation envelope as \(c_Y+\sum_ba_{b,Y}v_b^2\), with nonnegative coefficients uniformly bounded by Assumption \(\mathrm{ass:propensity\text{-}floor}\).  Hence its square root is the supremum of the norms
\(\alpha^{-1/2}\sqrt{c_Y+\sum_ba_{b,Y}v_b^2}\), and is convex and uniformly Lipschitz on bounded sets.  Definition \(\mathrm{def:drift\text{-}support\text{-}family}\) makes the bias term a support function of a bounded set, so it too is convex and Lipschitz.  Thus \(f\) is a finite continuous convex function on finite-dimensional space.  Since \(f(u)<2M\) and \(\min_Cf\le f(u)\), clipping is inactive at the optimum, and therefore
\[
\mathfrak R_{\alpha,\gamma}^\star=\min_{v\in C}f(v).
\]
It follows immediately that
\[
\mathfrak R_{\alpha,\gamma}^\star<f(u)
\quad\Longleftrightarrow\quad
u\notin\operatorname*{argmin}_{v\in C}f(v).
\]

For completeness, use the normal-cone convention
\[
N_C(u)=\{n:\langle n,v-u\rangle\le0\text{ for every }v\in C\}.
\]
If \(0=g+n\) with \(g\in\partial f(u)\) and \(n\in N_C(u)\), then for every \(v\in C\),
\[
f(v)\ge f(u)+\langle g,v-u\rangle
=f(u)-\langle n,v-u\rangle\ge f(u),
\]
so \(u\) minimizes \(f\) over \(C\).  For the converse, let
\[
A=\partial f(u),
\qquad
K=-N_C(u)=\{g:\langle g,h\rangle\ge0\text{ for every }h\in T_u\}.
\]
Because \(f\) is finite and continuous on finite-dimensional space, its secant inequalities give a nonempty compact convex subdifferential \(A\), and
\[
f'(u;h)=\max_{g\in A}\langle g,h\rangle.
\]
The latter identity follows directly by taking subgradients at \(u+t h\): local Lipschitz continuity bounds those subgradients; along any sequence \(t\downarrow0\), a convergent subsequence and the two subgradient inequalities at \(u\) and \(u+t h\) sandwich the difference quotient between the limiting inner products.

Suppose that \(u\) minimizes \(f\) on \(C\) but \(A\cap K=\varnothing\).  Choose \(a_0\in A\) and \(k_0\in K\) minimizing \(\|a-k\|\), and put \(w=a_0-k_0\ne0\).  Comparing the minimizing pair with \((a,k_0)\) and \((a_0,k)\), then taking one-sided derivatives of the squared distance, gives
\[
\langle w,a-a_0\rangle\ge0\quad(a\in A),
\qquad
\langle w,k-k_0\rangle\le0\quad(k\in K).
\]
Because \(K\) is a cone, the second inequality with \(k=0\) and \(k=2k_0\) gives \(\langle w,k_0\rangle=0\), and scaling any \(k\in K\) gives \(\langle w,k\rangle\le0\).  Hence, for \(h=-w\),
\[
\sup_{a\in A}\langle a,h\rangle
\le-\langle a_0,w\rangle
=-\|w\|^2<0,
\qquad
\langle k,h\rangle\ge0\quad(k\in K).
\]
The explicit form of \(K\) forces \(\sum_bh_b=0\) and \(h_b\ge0\) whenever \(u_b=0\): constants of both signs belong to \(K\), and so does every nonnegative coordinate vector off the support of \(u\).  Thus \(h\in T_u\).  The directional-derivative identity now gives \(f'(u;h)<0\), contradicting optimality along a sufficiently short feasible segment.  Therefore \(A\cap K\ne\varnothing\); taking \(g\in A\cap K\) gives \(-g\in N_C(u)\), and
\[
u\in\operatorname*{argmin}_{v\in C}f(v)
\quad\Longleftrightarrow\quad
0\in\partial f(u)+N_C(u).
\]

The cone of feasible one-sided directions at \(u\) is exactly
\[
T_u=\{h:\sum_bh_b=0,\ h_b\ge0\text{ whenever }u_b=0\}.
\]
To verify this, necessity follows from \(u+th\in C\) for arbitrarily small \(t>0\); conversely, for \(h\in T_u\), choosing
\(t>0\) below every positive ratio \(u_b/(-h_b)\) with \(h_b<0\) ensures \(u+th\in C\).  If \(u\) is not a minimizer, choose \(v\in C\) with \(f(v)<f(u)\) and let \(h=v-u\in T_u\).  Convexity of \(t\mapsto f(u+th)\) gives
\[
f'(u;h)\le f(v)-f(u)<0.
\]
Conversely, if \(h\in T_u\) and \(f'(u;h)<0\), the definition of the one-sided directional derivative gives \(f(u+th)<f(u)\) for all sufficiently small positive feasible \(t\).  Hence
\[
u\notin\operatorname*{argmin}_{v\in C}f(v)
\quad\Longleftrightarrow\quad
\exists h\in T_u:\ f'(u;h)<0.
\]
Combining the three equivalences proves the certificate.  The definitions of \(\mathcal G_{\mathrm{cc}}\) and \(\Gamma_{\mathrm{cc}}\) merely apply this certificate separately at each budget; no monotonicity of the schedule supremum in the stochastic term is used.  Finally, the optimization compares the full-record AIPW weight \(v_\gamma^\star\) with the concurrent AIPW weight \(u\); Definition \(\mathrm{def:concurrent\text{-}only\text{-}benchmark}\) separately defines the all-kernel value \(\mathfrak L_\alpha^{\mathrm{cc}}\).

%%% FIELD statement-strict-replacement
For each \(0\le\gamma\le\Gamma_{\mathrm{vac}}\) with \(\Psi_{\alpha,\gamma}(u)<2M\), the certified NCC AIPW radius \(\mathfrak R_{\alpha,\gamma}^\star\) is strictly shorter than the certified concurrent-only AIPW radius \(\Psi_{\alpha,\gamma}(u)\) if and only if \(u\notin\operatorname*{argmin}_{v\in\Delta(\mathcal B)}\Psi_{\alpha,\gamma}(v)\).  This is equivalent to \(0\notin\partial\Psi_{\alpha,\gamma}(u)+N_{\Delta(\mathcal B)}(u)\), and equivalent to the existence of \(h\in T_u\) with \(\Psi_{\alpha,\gamma}'(u;h)<0\).  The set \(\mathcal G_{\mathrm{cc}}\) and permanent threshold \(\Gamma_{\mathrm{cc}}\) are therefore computed from the same pointwise optimizer certificates without assuming that adaptive average-information terms are monotone in \(\gamma\).  This certificate compares the two AIPW constructions; the exact all-kernel concurrent benchmark is \(\mathfrak L_\alpha^{\mathrm{cc}}\).

%%% FIELD statement-mixture-lower
For every \(0\le\eta<1-2\alpha\), every schedule-independent kernel \(K\in\mathcal C_\alpha\) satisfies
\[
\sup_{Y\in\mathcal Y_{M,\Gamma}}
\mathbb E^K_{\mathcal D,Y}[\operatorname{len}(I_K)]
\ge (1-2\alpha-\eta)\Omega_{\mathcal D,\Gamma}(\eta).
\]
Consequently,
\[
\mathfrak L_\alpha^\star
\ge \eta_\alpha\Omega_{\mathcal D,\Gamma}(\eta_\alpha).
\]

%%% FIELD proof-mixture-lower
Fix \(K\in\mathcal C_\alpha\), \(0\le\eta<1-2\alpha\), and any admissible pair of priors \(\Pi_0,\Pi_1\) in Definition \(\mathrm{def:average\text{-}information\text{-}modulus}\) with target separation \(\Delta\) and mixture-law total variation at most \(\eta\).  Include the independent auxiliary randomizer in the data variable \(X=(O,U_{\mathrm{aux}})\), and denote its two mixture laws by \(Q_0,Q_1\).  Tensoring both laws with the same independent uniform distribution does not change their total variation, so \(\operatorname{TV}(Q_0,Q_1)\le\eta\).

For \(j\in\{0,1\}\), let \(a_j(x)\) be the conditional probability, under the \(j\)-th prior and its extended experiment, that the interval \(I_K(x)\) covers the corresponding random value \(\tau_W(Y)\).  Pointwise honesty and averaging over \(\Pi_j\) give
\[
\int a_j\,dQ_j\ge1-\alpha.
\]
Let \(\lambda=Q_0+Q_1\), let \(q_j\) be the density of \(Q_j\) with respect to \(\lambda\), and define the common subprobability measure \(\nu\) by
\(d\nu=\min(q_0,q_1)d\lambda\).  Directly from the definition of total variation,
\[
\nu(\mathsf X)=1-\operatorname{TV}(Q_0,Q_1),
\qquad \nu\le Q_j,
\qquad (Q_j-\nu)(\mathsf X)=\operatorname{TV}(Q_0,Q_1).
\]
If \(\operatorname{len}(I_K(x))<\Delta\), the support-separation condition prevents this interval from containing a target value from both prior supports.  Therefore
\[
\mathbf 1\{\operatorname{len}(I_K(x))\ge\Delta\}
\ge a_0(x)+a_1(x)-1.
\]
Since \(0\le a_j\le1\), domination by the common part gives
\[
\int a_j\,d\nu
\ge\int a_j\,dQ_j-(Q_j-\nu)(\mathsf X)
\ge1-\alpha-\operatorname{TV}(Q_0,Q_1).
\]
Integrating the indicator inequality against \(\nu\) yields
\[
\nu\{\operatorname{len}(I_K)\ge\Delta\}
\ge1-2\alpha-\operatorname{TV}(Q_0,Q_1)
\ge1-2\alpha-\eta.
\]
It follows that
\[
\begin{aligned}
\sup_{Y\in\mathcal Y_{M,\Gamma}}
 \mathbb E^K_{\mathcal D,Y}[\operatorname{len}(I_K)]
&\ge \int \mathbb E^K_{\mathcal D,Y}[\operatorname{len}(I_K)]\,d\Pi_0(Y)\\
&=\int\operatorname{len}(I_K(x))\,dQ_0(x)\\
&\ge\int\operatorname{len}(I_K(x))\,d\nu(x)\\
&\ge\Delta(1-2\alpha-\eta).
\end{aligned}
\]
Taking the supremum over all admissible prior pairs and separations, then the infimum over \(K\in\mathcal C_\alpha\), proves the first display.  At \(\eta=\eta_\alpha=(1-2\alpha)/2\), its multiplier is exactly \(\eta_\alpha\), proving the second display.  The proof covers variable-length kernels and auxiliary randomization because both are included in \(X\) and the same extended law is used in the coverage and length calculations.

%%% FIELD partial-oeq
The all-kernel testing step can be completed without any Gaussian or independent-sample approximation: the proved lemma \(\mathrm{lem:mixture\text{-}modulus\text{-}length\text{-}lower\text{-}bound}\) gives
\[
\mathfrak L_\alpha^\star\ge
\eta_\alpha\Omega_{\mathcal D,\Gamma}(\eta_\alpha).
\]
What remains open is the reverse comparison between the optimized AIPW radius and this observed-law modulus, namely the construction, for the arbitrary adaptive design in the frozen assumptions, of TV-admissible schedule-supported priors with separation of order \(\mathfrak R_{\alpha,\gamma}^\star\).

%%% FIELD partial-minimax
The constructive theorem proves the unconditional upper half with the explicit universal choice
\[
\mathfrak L_\alpha^\star\le2\mathfrak R_\alpha^\star.
\]
The proved mixture lemma gives the unconditional all-kernel lower bound
\[
\mathfrak L_\alpha^\star\ge
\eta_\alpha\Omega_{\mathcal D,\Gamma}(\eta_\alpha).
\]
Consequently, if the still-open comparison
\(\mathfrak R_\alpha^\star\le C_\alpha
\Omega_{\mathcal D,\Gamma}(\eta_\alpha)\) is established, then the claimed lower half follows with
\(c_\alpha^{\mathrm{lb}}=\eta_\alpha/C_\alpha\), while one may take
\(C_\alpha^{\mathrm{ub}}=2\).  The frozen primitives do not presently prove the required comparison.

%%% FIELD partial-boundary
Several boundary assertions follow directly from the frozen core.  First, at \(\gamma=0\), every feasible control-mean vector is constant; since \(\sum_bd_b(v)=0\), Lemma \(\mathrm{lem:tv\text{-}support\text{-}dual}\) gives
\[
\mathsf B_0(v)=0.
\]
Thus \(v_0^\star\) minimizes only the stochastic radius over the full simplex.  Second, for \(\gamma\in\mathcal G_{\mathrm{cc}}\), the KKT equivalence proved in Proposition \(\mathrm{prop:strict\text{-}improvement\text{-}iff}\) says exactly that \(u\) is optimal.  Third, Definition \(\mathrm{def:concurrent\text{-}only\text{-}benchmark}\) gives \(\mathcal C_\alpha^{\mathrm{cc}}\subseteq\mathcal C_\alpha\), and therefore
\[
\mathfrak L_\alpha^\star\le\mathfrak L_\alpha^{\mathrm{cc}}.
\]
Finally, substituting \(v=u\) in Definition \(\mathrm{def:borrowing\text{-}estimator}\) sets every off-window control coefficient to zero and leaves precisely the window-weighted treatment-minus-control AIPW contrast.

Two clauses remain uncertified.  Strict convexity by itself does not force a simplex minimizer to be interior; the positive-weight claim needs an explicitly quantified uniform lower bound on every block's quadratic coefficient, not merely pointwise positive design probability.  Also, the reverse inequality at \(\Gamma=\Gamma_{\mathrm{vac}}\) does not follow from the stated reduction.  Under history-adaptive randomization, the conditional law of the omitted off-window outcomes given \(O^{\mathrm{cc}}\) can depend on the unknown fixed schedule.  Averaging a full-record kernel against such a conditional law can therefore produce a schedule-dependent kernel, which is not an element of \(\mathcal C_\alpha^{\mathrm{cc}}\).
